\chapter{Аналитический обзор}
\label{ch:chap1}

\section{Имитационное обучение для задач манипулирования}

Имитационное обучение формирует политику агента на основе экспертных демонстраций без определения
функции вознаграждения. Простейшая форма, клонирование поведения, сводится к регрессии действий по
наблюдениям. Мандлекар и соавторы \cite{mandlekar2021robomimic} показали, что человеческие демонстрации
порождаются немарковским процессом с мультимодальным распределением действий, из-за чего усреднение по
режимам приводит к некорректной политике.

Для преодоления этих ограничений предложены несколько подходов. BC-RNN учитывает временные зависимости
через скрытое состояние и существенно превосходит стандартное клонирование поведения
\cite{mandlekar2021robomimic}. Диффузионные политики \cite{chi2023diffusion} моделируют
мультимодальность через итеративное шумоподавление, но требуют многократных проходов при выводе.
ACT \cite{zhao2023aloha} объединяет трансформер с предсказанием последовательностей действий.
Behavior Transformer \cite{shafiullah2022behavior} дискретизирует действия кластеризацией $k$-means,
однако этот метод плохо масштабируется на высокоразмерные пространства, что мотивирует применение
обучаемых методов дискретизации.

\section{Визуомоторные политики и проблема робастности}

Визуомоторная политика преобразует изображения в команды актуаторам через визуальный энкодер
(обычно ResNet) и генератор действий. По данным \cite{mandlekar2021robomimic}, BC-RNN с аугментацией
изображений практически не уступает агентам с доступом к низкоразмерным состояниям; при этом удаление
пиксельных сдвигов снижает успешность на 35–47\%.

Доменные сдвиги -- изменения освещения, фона, текстур, появление окклюзий -- остаются ключевым
препятствием для развёртывания. Росс и соавторы \cite{ross2011reduction} показали, что ошибки политики
накапливаются квадратично по длине горизонта (covariate shift), а визуальные сдвиги усугубляют этот
эффект. Стандартные меры -- аугментация и доменная рандомизация -- действуют на уровне входных данных.
Ян и соавторы \cite{yang2023vector} показали, что модуль векторного квантования в визуальном энкодере
создаёт дискретное «бутылочное горлышко», подавляющее нерелевантные компоненты и повышающее
инвариантность представлений.

\section{Квантование представлений и его применение в робототехнике}

Векторное квантование заменяет непрерывный вектор ближайшим прототипом из конечной кодовой книги.
Ван ден Оорд и соавторы предложили VQ-VAE \cite{vandeoord2017neural}, в котором дискретные латентные
переменные устраняют коллапс постериора и формируют информационное «бутылочное горлышко». Зегиду и
соавторы расширили подход до остаточного квантования (Residual VQ) \cite{zeghidour2021soundstream}:
каскад слоёв обеспечивает гибкий компромисс между компактностью и точностью.

В робототехнике квантование применяется по нескольким направлениям. Квантование действий: VQ-BeT
\cite{lee2024vqbet} заменяет $k$-means в BeT обучаемым Residual VQ-VAE; на семи средах метод превзошёл
диффузионные политики при пятикратном ускорении вывода, а на длинногоризонтных задачах разрыв составил
73\%. Куи и соавторы \cite{cui2026contact} использовали VQ-BeT для политик, обусловленных точками
контакта, и превзошли крупные VLA-модели на 56\%. Квантование состояний: Фэн и соавторы
\cite{feng2025learning} квантовали состояния кисти двухслойным Residual VQ-VAE, преобразовав
предсказание высокоразмерных действий в задачу классификации. Квантование визуальных наблюдений:
Маджан и соавторы \cite{mahajan2020robotic} применили VQ-VAE для обнаружения захватов, улучшив обобщение
при ограниченных данных; Порихис и соавторы \cite{porichis2024imitation} достигли 100\% успешности
захвата на единственной демонстрации.

\section{Среда моделирования и набор данных}

Бенчмарк robomimic \cite{mandlekar2021robomimic} включает пять симулированных и три реальных задачи
возрастающей сложности с демонстрациями от операторов различной квалификации. Среда robosuite на основе
MuJoCo обеспечивает реалистичную симуляцию и поддерживает программное изменение освещения, текстур и
положения камеры, что позволяет систематически исследовать робастность к доменным сдвигам. В настоящей
работе используются robomimic v0.1 и robosuite в качестве основной экспериментальной платформы.

\section{Выводы по результатам обзора}

Векторное квантование продемонстрировало эффективность при дискретизации действий, состояний и визуальных
наблюдений в робототехнике. Вместе с тем систематического исследования влияния квантования визуальных
представлений на обучаемость и робастность в задаче манипулирования не проводилось. Данный пробел
определяет направление работы: оценка способности модуля VQ в архитектуре BC-RNN одновременно повысить
обучаемость и устойчивость к доменным сдвигам при захвате объектов.

\endinput
